\section{Threat Model and Taxonomy}\label{sec:threats}
\answers{RQ1}

\paragraph*{Adversaries} (A1)~A \emph{content author} controls text the agent reads (README, comment, issue, web page, tool result); (A2)~a \emph{repository author} controls cloned files, including configuration; (A3)~a \emph{supply-chain attacker} controls a package, MCP server, plugin or release; (A4)~a \emph{local or network peer} reaches a service the harness exposes. The defender runs the harness with real credentials, code and host access (the assets). We assume the model sometimes follows injected instructions and that no model-level defence holds.

\paragraph*{Taxonomy} We extended the 16-class seed into the classes in Table~\ref{tab:taxonomy}. Three classes were added because the incident set contained large groups with no home: sandbox and isolation escape (T17), exposure of the agent's own control plane, such as local HTTP or WebSocket servers without authentication or Origin checks (T18), and conventional software defects in runtimes and SDKs (T19). Each class carries STRIDE letters, MITRE ATLAS techniques, OWASP LLM~2025, OWASP Agentic (ASI01--ASI10) and OWASP MCP Top~10 identifiers; the MCP list is still in beta, and we flag those mappings as unstable. Each class also lists the control dimensions that mitigate it (Section~\ref{sec:controls}) and the \tool{} rules that detect it (Appendix~\ref{app:rules}).

\paragraph*{Classification procedure} Write-ups mix the entry vector with the control that failed, so we walk the classes in a fixed order and assign the first whose \emph{failed control} the incident exhibits. A parallel 12-question procedure assigns one of twelve root causes, RC1--RC12 (Table~\ref{tab:rootcauses}). For example, injection-written configuration that later executes is T02/RC2, not T01; data leaving through a permitted command is T08/RC3. RC1 (untrusted content treated as instructions with no other control failure) is the residual answer reached only when no earlier question applies, so its share is low partly by construction.

% generated by experiments/taxonomy_table.py; do not edit
\begin{table*}[t]
\centering\scriptsize\setlength{\tabcolsep}{3pt}
\caption{Harness threat taxonomy (RQ1) with framework mappings, mitigating control dimensions (D01--D15, Section~\ref{sec:controls}) and the number of verified 2025--2026 incidents whose primary class it is (RQ3). ATLAS ids are abbreviated (prefix \texttt{AML.} dropped; first shown).}
\label{tab:taxonomy}
\begin{tabular}{lp{5.3cm}lp{1.9cm}p{2.3cm}p{1.7cm}p{2.9cm}r}
\toprule
Class & Name & STRIDE & OWASP LLM & OWASP ASI & MITRE ATLAS & Controls & Inc. \\
\midrule
T01 & Indirect prompt injection via workspace and retrieved content & TE & LLM01 & ASI01 & T0051.001, \dots & D01, D02, D03, D04, D11 & 14 \\
T02 & Repository and workspace configuration as code execution & ET & LLM06 & ASI05, ASI02 & T0081, \dots & D04, D01, D02, D05, D14 & 52 \\
T03 & Permission-model weaknesses and approval bypass & ET & LLM06 & ASI02, ASI03 & T0053, \dots & D01, D02, D14 & 74 \\
T04 & MCP trust: tool poisoning, rug pulls, OAuth and server identity & STE & LLM03, LLM01, LLM06 & ASI04, ASI02 & T0110, \dots & D05, D14, D06, D13, D03 & 15 \\
T05 & Skills, plugins, extensions and marketplaces & STE & LLM03 & ASI04 & T0010.005, \dots & D06, D13, D14 & 2 \\
T06 & Harness distribution and dependency supply chain & TSE & LLM03 & ASI04 & T0010.001, \dots & D13, D15, D14 & 9 \\
T07 & Secret exposure and residue & I & LLM02 & ASI03 & T0055, \dots & D07, D11, D12 & 12 \\
T08 & Data exfiltration channels & I & LLM02, LLM05 & ASI02 & T0086, \dots & D03, D02, D01 & 23 \\
T09 & Headless and CI abuse (pwn request) & EIT & LLM01, LLM06 & ASI01, ASI02, ASI03 & T0051.001, \dots & D08, D01, D14, D07 & 5 \\
T10 & Persistent memory and long-lived context poisoning & TE & LLM04 & ASI06 & T0080.000, \dots & D09, D11 & 2 \\
T11 & Sub-agent and cross-harness trust propagation & ES & LLM06 & ASI03, ASI07, ASI08 & T0118, \dots & D10, D01, D02 & 2 \\
T12 & Browser and computer-use surfaces & SIE & LLM01, LLM06 & ASI01, ASI02 & T0100, \dots & D02, D03, D01 & 1 \\
T13 & Multi-tenant isolation and authorization & IES & LLM02 & ASI03 & T0085, \dots & D14, D10, D07 & 7 \\
T14 & Destructive and unsafe autonomous actions & DT & LLM06 & ASI02, ASI08 & T0101, \dots & D01, D02, D14 & 0 \\
T15 & Telemetry, data residency, retention and audit gaps & IR & LLM02 &  & -- & D12, D11, D14 & 1 \\
T16 & Offensive misuse of harnesses by attackers & E &  & ASI10 & T0124, \dots & D11, D14 & 3 \\
T17 & Execution sandbox and isolation escape & E & LLM06 & ASI05 & T0105, \dots & D02, D10, D03 & 22 \\
T18 & Agent control-plane exposure and authorization & SEI &  & ASI03 & T0132, \dots & D16, D01, D04 & 34 \\
T19 & Conventional software defects in agent runtimes and SDKs & TDE & LLM05, LLM10 & ASI05 & T0049, \dots & D15, D13, D02 & 45 \\
\bottomrule
\end{tabular}
\end{table*}


\paragraph*{How harness risk differs from LLM-application risk} In a chat application, a successful injection yields a bad answer or a leaked system prompt. In a harness, the injected text comes from the developer's own workspace and is acted on with the developer's privileges. Three properties follow. First, \emph{configuration is code}: repository files can define hooks, environment variables and MCP launch commands that run before the model is consulted (T02). Second, \emph{the approval layer is a parser}: allowlists over shell commands must reason about quoting, substitution, redirection and symlinks, and such failures are the largest single group of disclosed incidents (T03, \pctTopThree{} by class; RC3, \pctRCThree{} by root cause). Third, \emph{the harness is a long-lived service}: local servers, sub-agents, memory and CI runners create surfaces that classic LLM threat models do not include (T10, T11, T18). Attack trees for the main paths are in the repository.
